%%======================================================================
%%  NT-4c: Literature Extraction & Convention Fixing
%%         for FLRW Reduction of Nonlocal Field Equations
%%
%%  This document collects formulas from the literature needed to
%%  reduce the NT-4b nonlinear field equations to FLRW symmetry,
%%  derive modified Friedmann equations, and extract the perturbation
%%  framework for tensor and scalar modes.
%%
%%  RESEARCH DOCUMENT: formulas extracted from cited papers.
%%  No original derivations --- gaps noted explicitly.
%%======================================================================
\documentclass[12pt,a4paper]{article}

%% ---- packages -------------------------------------------------------
\usepackage[utf8]{inputenc}
\usepackage[T1]{fontenc}
\usepackage{lmodern}
\usepackage{amsmath,amssymb,amsthm,mathtools}
\usepackage[margin=2.5cm]{geometry}
\usepackage{booktabs}
\usepackage{enumitem}
\usepackage{hyperref}
\usepackage{xcolor}
\usepackage{fancyhdr}
\usepackage{titlesec}
\usepackage{longtable}
\usepackage{tcolorbox}

%% ---- theorem-like environments --------------------------------------
\newtheorem{proposition}{Proposition}[section]
\newtheorem{lemma}[proposition]{Lemma}
\newtheorem{corollary}[proposition]{Corollary}
\newtheorem{theorem}[proposition]{Theorem}
\theoremstyle{definition}
\newtheorem{definition}[proposition]{Definition}
\newtheorem{convention}[proposition]{Convention}
\theoremstyle{remark}
\newtheorem{remark}[proposition]{Remark}

%% ---- custom commands ------------------------------------------------
\newcommand{\Tr}{\operatorname{Tr}}
\newcommand{\tr}{\operatorname{tr}}
\newcommand{\Rsq}{R_{\mu\nu\rho\sigma}R^{\mu\nu\rho\sigma}}
\newcommand{\Ricsq}{R_{\mu\nu}R^{\mu\nu}}
\newcommand{\Weylsq}{C_{\mu\nu\rho\sigma}C^{\mu\nu\rho\sigma}}
\newcommand{\dd}{\mathrm{d}}
\newcommand{\ii}{\mathrm{i}}
\newcommand{\Id}{\mathrm{Id}}
\newcommand{\erfi}{\operatorname{erfi}}
\DeclareMathOperator{\sgn}{sgn}

%% ---- annotation commands --------------------------------------------
\newcommand{\fromlit}[1]{\textcolor{blue!70!black}{\textsf{[#1]}}}
\newcommand{\oursign}{\textcolor{teal!80!black}{\textsf{[our convention]}}}
\newcommand{\gap}{\textcolor{red!80!black}{\textsf{[GAP]}}}
\newcommand{\needscheck}{\textcolor{orange!80!black}{\textsf{[needs check]}}}
\newcommand{\exactlit}{\textcolor{blue!60!black}{\textsf{[exact from lit.]}}}
\newcommand{\verified}[1]{%
  \par\smallskip\noindent
  \colorbox{green!10}{\parbox{\dimexpr\linewidth-2\fboxsep}{%
    \textcolor{green!50!black}{\textbf{[Verified]}} #1}}%
  \par\smallskip}
\newcommand{\signcorrection}[1]{%
  \par\noindent\colorbox{red!10}{%
    \parbox{\dimexpr\linewidth-2\fboxsep}{%
      \textcolor{red!50!black}{\textbf{[Sign correction]}}
      \textcolor{red!50!black}{#1}%
    }%
  }\par%
}

%% ---- page style -----------------------------------------------------
\pagestyle{fancy}
\fancyhf{}
\fancyhead[L]{\small\textsc{NT-4c: Literature \& Conventions for
FLRW Reduction}}
\fancyhead[R]{\small\thepage}
\renewcommand{\headrulewidth}{0.4pt}

%% ---- hyperref -------------------------------------------------------
\hypersetup{
  colorlinks=true,
  linkcolor=blue!60!black,
  citecolor=green!40!black,
  urlcolor=blue!60!black,
  pdftitle={NT-4c: Literature Extraction for FLRW Reduction of Nonlocal Field Equations},
  pdfauthor={David Alfyorov},
}

%% ---- title ----------------------------------------------------------
\title{\textbf{NT-4c: Literature Extraction \& Convention Fixing\\[4pt]
       for FLRW Reduction of Nonlocal Field Equations}\\[8pt]
       \large From Nonlinear SCT Field Equations to Modified Friedmann
       Equations\\and Cosmological Perturbation Framework}
\author{David Alfyorov}
\date{March 2026 --- Working Document}

%%======================================================================
\begin{document}
\maketitle

\begin{center}
\small\textit{Credits: Aliaksandr Samatyia contributed research-assistance
and workflow support.}
\end{center}

\thispagestyle{empty}

\begin{abstract}
This document collects all formulas, conventions, and computational
recipes from the literature needed to reduce the full nonlinear SCT
field equations (NT-4b) to spatially flat FLRW symmetry and derive
the modified Friedmann equations. The central simplification is that
the Weyl tensor vanishes identically on conformally flat backgrounds
($C_{\mu\nu\rho\sigma} = 0$ on FLRW), so the entire $C^2$ sector
governed by $F_1(\Box/\Lambda^2)$ drops out of the \emph{background}
equations. Only the $R^2$ sector, controlled by
$\alpha_R(\xi) = 2(\xi - 1/6)^2$, contributes nonlocal corrections
to the Friedmann equations. At conformal coupling $\xi = 1/6$,
$\alpha_R = 0$ and the spectral action produces no background
modifications beyond the Einstein--Hilbert term.

The document covers: (1)~FLRW geometry and curvature identities;
(2)~the d'Alembertian on FLRW and its action on scalars;
(3)~reduction of the $R^2$ sector including the
$\Theta^{(R)}_{\mu\nu}$ alpha-insertion terms;
(4)~the resulting modified Friedmann equations;
(5)~de~Sitter self-consistency;
(6)~scalar, vector, and tensor perturbations on FLRW;
(7)~gravitational wave speed and the GW170817 constraint;
(8)~comparison with Deser--Woodard, Maggiore--Mancarella,
Koshelev--Kumar--Starobinsky, and Biswas--Koshelev--Mazumdar
nonlocal cosmological models.
All formulas are extracted from the literature with source
annotations; no original derivations are performed. Gaps requiring
resolution in the derivation phase are noted explicitly.

\medskip\noindent
\textbf{Key references:}
\begin{itemize}[nosep]
\item NT-4b: Full nonlinear field equations (this project).
\item NT-4a: Linearized field equations (this project).
\item Stelle: Phys.\ Rev.\ D \textbf{16} (1977) 953;
  Gen.\ Rel.\ Grav.\ \textbf{9} (1978) 353.
\item Barvinsky--Vilkovisky (BV): Phys.\ Rep.\ \textbf{119} (1985) 1;
  Nucl.\ Phys.\ B \textbf{282} (1987) 163;
  B \textbf{333} (1990) 471.
\item Deser--Woodard (DW): Phys.\ Rev.\ Lett.\ \textbf{99} (2007) 111301
  \cite{DW07}; arXiv:1307.6639 \cite{DW13}; arXiv:1902.08075 \cite{DW19}.
\item Belgacem--Dirian--Foffa--Maggiore (BDFM): arXiv:1712.07066
  \cite{BDFM18}.
\item Dirian--Foffa--Khosravi--Kunz--Maggiore (DFKKM):
  arXiv:1403.6068 \cite{DFKKM14}.
\item Koshelev--Kumar--Starobinsky (KKS): arXiv:2305.18716 \cite{KKS23}.
\item Biswas--Mazumdar--Siegel (BMS): hep-th/0508194 \cite{BMS06}.
\item Biswas--Koivisto--Mazumdar (BKM): arXiv:1005.0590 \cite{BKM10}.
\item Briscese--Calcagni--Modesto (BCM): arXiv:1901.03267 \cite{BCM19}.
\item Calcagni--Modesto (CM): arXiv:2211.05606 \cite{CM22}.
\item Nersisyan--Lima--Amendola (NLA): arXiv:1801.06683 \cite{NLA18}.
\item Belgacem--Dirian--Foffa--Maggiore: arXiv:1805.08731 \cite{BDFM18gw}.
\item Nojiri--Odintsov (NO): arXiv:0810.1557 \cite{NO08}.
\end{itemize}
\end{abstract}

\tableofcontents
\newpage

%%======================================================================
\section{Input: The NT-4b Field Equations}\label{sec:input}
%%======================================================================

\subsection{The Canonical Field Equations (NT-4b, Theorem 4.1)}

The starting point for the FLRW reduction is the full nonlinear
field equation derived in NT-4b:
\begin{equation}\label{eq:full-eom}
  \boxed{
  \frac{1}{\kappa^2}\,G_{\mu\nu}
  + \frac{1}{16\pi^2}\bigl[2\,F_1(\Box/\Lambda^2)\,B_{\mu\nu}
    + \Theta^{(C)}_{\mu\nu}\bigr]
  + \frac{1}{16\pi^2}\bigl[F_2(\Box/\Lambda^2,\xi)\,H_{\mu\nu}
    + \Theta^{(R)}_{\mu\nu}\bigr]
  = \frac{1}{2}\,T_{\mu\nu},
  }
\end{equation}
where \fromlit{NT-4b, Eq.~(full-eom)}:
\begin{itemize}[nosep]
  \item $G_{\mu\nu} = R_{\mu\nu} - \tfrac{1}{2}g_{\mu\nu}R$
    is the Einstein tensor,
  \item $B_{\mu\nu} = \nabla^\rho\nabla^\sigma C_{\mu\rho\nu\sigma}
    + \tfrac{1}{2}R^{\rho\sigma}C_{\mu\rho\nu\sigma}$ is the Bach
    tensor (traceless, divergence-free),
  \item $H_{\mu\nu} = 2\nabla_\mu\nabla_\nu R - 2g_{\mu\nu}\Box R
    - \tfrac{1}{2}g_{\mu\nu}R^2 + 2R\,R_{\mu\nu}$ is the $R^2$-sector
    tensor with trace $g^{\mu\nu}H_{\mu\nu} = -6\Box R$,
  \item $\Theta^{(R)}_{\mu\nu}$ and $\Theta^{(C)}_{\mu\nu}$ are the
    BV alpha-insertion corrections,
  \item $\kappa^2 = 8\pi G$, $\Lambda$ is the spectral cutoff.
\end{itemize}

\subsection{Canonical Coefficients}

From NT-1b Phase~3 \fromlit{NT-1b, Phase~3 handoff}:
\begin{align}
  \alpha_C &= 16\pi^2\,F_1(0) = \frac{13}{120},
  \label{eq:alphaC}\\
  \alpha_R(\xi) &= 16\pi^2\,F_2(0,\xi) = 2\Bigl(\xi - \frac{1}{6}\Bigr)^2,
  \label{eq:alphaR}\\
  c_2 &= 2\alpha_C = \frac{13}{60},
  \label{eq:c2}\\
  3c_1 + c_2 &= 6\Bigl(\xi - \frac{1}{6}\Bigr)^2.
  \label{eq:3c1c2}
\end{align}

\subsection{The $\Theta^{(R)}$ Alpha-Insertion (NT-4b, \S4)}

\begin{equation}\label{eq:Theta-R}
  \Theta^{(R)}_{\mu\nu} = -\sum_{n=1}^{\infty} c_n^{(2)}
  \sum_{k=0}^{n-1}
  \Bigl[\nabla_\mu(\Box^k R)\,\nabla_\nu(\Box^{n-1-k} R)
  - \frac{1}{2}g_{\mu\nu}\,\nabla_\rho(\Box^k R)\,
  \nabla^\rho(\Box^{n-1-k} R)\Bigr],
\end{equation}
where $F_2(z,\xi) = \sum_{n=0}^{\infty} c_n^{(2)}\,(z/\Lambda^2)^n$.

\verified{NT-4b: 110/110 tests pass. Dual derivation (Methods A and
B1) exact agreement. De Sitter, Bianchi, trace, local limit all
verified.}

%%======================================================================
\section{FLRW Geometry}\label{sec:flrw}
%%======================================================================

\subsection{Metric and Symmetry}

The spatially flat FLRW metric in cosmic time is \fromlit{Weinberg
(2008), \S1.1; Mukhanov (2005), \S5.1}:
\begin{equation}\label{eq:flrw-metric}
  \dd s^2 = -\dd t^2 + a(t)^2\,\delta_{ij}\,\dd x^i\,\dd x^j,
\end{equation}
with scale factor $a(t)$, Hubble parameter $H \equiv \dot{a}/a$,
and spatial curvature $k = 0$.

In conformal time $\tau$ defined by $\dd\tau = \dd t/a(t)$:
\begin{equation}\label{eq:flrw-conformal}
  \dd s^2 = a(\tau)^2\bigl(-\dd\tau^2 + \delta_{ij}\,\dd x^i\,\dd x^j\bigr).
\end{equation}

\begin{convention}[Signature and time]
We use Lorentzian signature $(-,+,+,+)$ throughout this document,
consistent with MR-1 (Wick rotation from Euclidean). Cosmic time
derivatives are denoted by overdots, conformal time derivatives by
primes.
\end{convention}

\subsection{Curvature Tensors on FLRW}\label{sec:flrw-curvature}

The nonvanishing components of the Ricci tensor and Ricci scalar on
the metric~\eqref{eq:flrw-metric} are
\fromlit{Weinberg (2008), \S1.5; Carroll (2004), \S8.3}:
\begin{align}
  R_{00} &= -3(\dot{H} + H^2) = -3\frac{\ddot{a}}{a},
  \label{eq:R00}\\
  R_{ij} &= \bigl(\dot{H} + 3H^2\bigr)\,a^2\,\delta_{ij}
  = \bigl(a\ddot{a} + 2\dot{a}^2\bigr)\,\delta_{ij},
  \label{eq:Rij}\\
  R &= 6\bigl(\dot{H} + 2H^2\bigr)
  = 6\Bigl(\frac{\ddot{a}}{a} + H^2\Bigr).
  \label{eq:R-flrw}
\end{align}

The Riemann tensor on a spatially flat FLRW metric with $k = 0$
has components \fromlit{Carroll (2004), \S8.3; Weinberg (2008), \S1.5}:
\begin{equation}\label{eq:Riem-flrw}
  R^0{}_{i0j} = \frac{\ddot{a}}{a}\,\delta_{ij},
  \qquad
  R^i{}_{jkl} = H^2\bigl(\delta^i_k\,\delta_{jl}
  - \delta^i_l\,\delta_{jk}\bigr),
\end{equation}
where the spatial $\delta_{ij}$ is the coordinate Kronecker delta.
Note: in Cartesian FLRW coordinates $\dd s^2 = -\dd t^2 +
a^2\delta_{ij}\dd x^i\dd x^j$, the component with all spatial
indices lowered is $R_{0i0j} = g_{00}\,R^0{}_{i0j} = -\ddot{a}\,a
\,\delta_{ij}$ (includes an extra factor of $a^2$ relative to the
mixed component).

\begin{proposition}[Weyl tensor on FLRW]\label{prop:weyl-flrw}
The Weyl tensor vanishes identically on any FLRW spacetime:
\begin{equation}\label{eq:weyl-zero}
  C_{\mu\nu\rho\sigma}\big|_{\mathrm{FLRW}} = 0.
\end{equation}
\end{proposition}

\begin{proof}
An FLRW spacetime is conformally flat (the metric is conformal to
Minkowski space). The Weyl tensor is conformally invariant in $d = 4$
and vanishes on flat space, hence
$C_{\mu\nu\rho\sigma}|_{\mathrm{FLRW}} = 0$.
Alternatively, on a spacetime of the form $g_{\mu\nu} = a^2(\tau)\,
\eta_{\mu\nu}$, the Weyl tensor vanishes because the Riemann tensor
has the maximally symmetric structure.
\end{proof}

\begin{tcolorbox}[colback=blue!5!white, colframe=blue!50!black,
  title={\textbf{Critical Consequence: Weyl Sector Drops Out of
  Background}}]
Since $C_{\mu\nu\rho\sigma} = 0$ on FLRW:
\begin{enumerate}[nosep]
  \item The Bach tensor vanishes: $B_{\mu\nu}\big|_{\mathrm{FLRW}} = 0$.
  \item The entire $F_1(\Box/\Lambda^2)$ sector does not contribute
    to the \textbf{background} Friedmann equations.
  \item The alpha-insertion $\Theta^{(C)}_{\mu\nu}\big|_{\mathrm{FLRW}} = 0$
    since all terms involve $C_{\mu\nu\rho\sigma}$ or its derivatives.
  \item Only the $R^2$ sector survives, governed by
    $\alpha_R(\xi) = 2(\xi - 1/6)^2$.
  \item At conformal coupling $\xi = 1/6$: $\alpha_R = 0$, so all
    spectral corrections vanish and one recovers the standard
    Friedmann equations exactly.
\end{enumerate}
\end{tcolorbox}

\subsection{The Ricci Tensor in Maximally Symmetric Form}

On FLRW, the Ricci tensor takes the ``perfect fluid'' form:
\begin{equation}\label{eq:Ric-pf}
  R_{\mu\nu} = (\dot{H} + 3H^2)\,g_{\mu\nu}
  - 2\dot{H}\,u_\mu u_\nu,
\end{equation}
where $u^\mu = (1, 0, 0, 0)$ is the comoving 4-velocity ($u_\mu =
(-1, 0, 0, 0)$). One verifies: $R_{00} = -(\dot{H} + 3H^2) -
2\dot{H}\cdot 1 = -3(\dot{H} + H^2)$; $R_{ij} = (\dot{H} + 3H^2)\,
a^2\delta_{ij}$.

\signcorrection{LR audit correction: The original formula
$R_{\mu\nu} = (R/4)\,g_{\mu\nu} + (R/4 + 3\dot{H})\,u_\mu u_\nu$
is incorrect on general FLRW. It holds only on de~Sitter where
$\dot{H} = 0$ and $R/4 = 3H^2 = \dot{H} + 3H^2$. On general FLRW,
$R/4 = \frac{3}{2}(\dot{H} + 2H^2) \neq \dot{H} + 3H^2$.}

On exact de~Sitter ($\dot{H} = 0$, $R = 12H^2$):
\begin{equation}\label{eq:Ric-dS}
  R_{\mu\nu}\big|_{\mathrm{dS}} = \frac{R}{4}\,g_{\mu\nu} = 3H^2\,g_{\mu\nu}.
\end{equation}

%%======================================================================
\section{The d'Alembertian on FLRW}\label{sec:box-flrw}
%%======================================================================

\subsection{Scalar d'Alembertian}

For a scalar field $\phi(t)$ depending only on cosmic time
\fromlit{Weinberg (2008), \S10.8; Mukhanov (2005), \S5.2}:
\begin{equation}\label{eq:box-scalar}
  \Box\phi = g^{\mu\nu}\nabla_\mu\nabla_\nu\phi
  = -\ddot{\phi} - 3H\dot{\phi}.
\end{equation}

In conformal time:
\begin{equation}\label{eq:box-scalar-conformal}
  \Box\phi = -\frac{1}{a^2}\bigl(\phi'' + 2\mathcal{H}\phi'\bigr),
\end{equation}
where $\mathcal{H} = a'/a = aH$ is the conformal Hubble parameter.

\subsection{$\Box R$ on FLRW}\label{sec:box-R}

Since $R = R(t)$ on FLRW, applying~\eqref{eq:box-scalar}:
\begin{equation}\label{eq:boxR}
  \Box R = -\ddot{R} - 3H\dot{R}.
\end{equation}

The explicit expression requires differentiating~\eqref{eq:R-flrw}:
\begin{align}
  \dot{R} &= 6(\ddot{H} + 4H\dot{H}),
  \label{eq:Rdot}\\
  \ddot{R} &= 6\bigl(\dddot{H} + 4\dot{H}^2 + 4H\ddot{H}\bigr).
  \label{eq:Rddot}
\end{align}

\begin{remark}[De Sitter limit]
On de~Sitter ($H = \mathrm{const}$): $\dot{R} = 0$,
$\Box R = 0$, and $\Box^n R = 0$ for all $n \geq 1$.
\end{remark}

\subsection{Iterated $\Box^n R$}\label{sec:boxnR}

For the alpha-insertion terms~\eqref{eq:Theta-R}, one needs
$\Box^n R$ for $n = 1, 2, \ldots$. On FLRW each $\Box^n R$ is a
function of $t$ alone, computed iteratively:
\begin{equation}\label{eq:boxnR-iter}
  \Box^{n+1} R = -\frac{\dd^2}{\dd t^2}(\Box^n R)
  - 3H\,\frac{\dd}{\dd t}(\Box^n R).
\end{equation}

\gap{} The convergence of the infinite series in~\eqref{eq:Theta-R}
on a general FLRW background must be established. The entireness of
$F_2$ (NT-2) ensures $|c_n^{(2)}| \sim 1/n!$, which suggests
convergence, but an explicit proof on FLRW is needed.

\subsection{$F_2(\Box/\Lambda^2)$ Acting on $R$}

The form factor acts on the Ricci scalar as:
\begin{equation}\label{eq:F2-action-R}
  F_2(\Box/\Lambda^2,\xi)\,R = \sum_{n=0}^{\infty}
  c_n^{(2)}\,\frac{\Box^n R}{\Lambda^{2n}}.
\end{equation}

On de~Sitter, this reduces to
$F_2(\Box/\Lambda^2,\xi)\,R = c_0^{(2)}\,R_0 = F_2(0,\xi)\,R_0$,
since $\Box^n R = 0$ for $n \geq 1$.

%%======================================================================
\section{Reduction of the $R^2$ Sector to FLRW}\label{sec:r2-flrw}
%%======================================================================

\subsection{The $H_{\mu\nu}$ Tensor on FLRW}

The $R^2$-sector tensor (NT-4b, Eq.~(H-tensor)) is:
\begin{equation}\label{eq:H-def}
  H_{\mu\nu} = 2\nabla_\mu\nabla_\nu R
  - 2g_{\mu\nu}\Box R
  - \tfrac{1}{2}g_{\mu\nu}R^2
  + 2R\,R_{\mu\nu}.
\end{equation}

On FLRW with $R = R(t)$, each term evaluates as follows.

\paragraph{(i) $\nabla_\mu\nabla_\nu R$:}
Since $R = R(t)$, $\nabla_\mu R = \dot{R}\,\delta^0_\mu$
\fromlit{standard: $\nabla_\mu R = \partial_\mu R$ for a scalar},
and:
\begin{align}
  \nabla_0\nabla_0 R &= \ddot{R},
  \label{eq:nnR-00}\\
  \nabla_i\nabla_j R &= -H\dot{R}\,a^2\,\delta_{ij}
  = -H\dot{R}\,g_{ij},
  \label{eq:nnR-ij}\\
  \nabla_0\nabla_i R &= 0.
  \label{eq:nnR-0i}
\end{align}

\signcorrection{LR audit correction:
$\nabla_0 R = +\dot{R}$ (positive), not $-\dot{R}$.
The raised-index gradient $\nabla^0 R = g^{00}\nabla_0 R = -\dot{R}$
has the minus sign; the lowered-index gradient does not.
$\nabla_0\nabla_0 R = +\ddot{R}$ (via
$\partial_0(\dot{R}) - \Gamma^\lambda_{00}\,\nabla_\lambda R
= \ddot{R} - 0$).
$\nabla_i\nabla_j R = -\Gamma^0_{ij}\,\nabla_0 R
= -(a\dot{a}\,\delta_{ij})\dot{R} = -H\dot{R}\,a^2\delta_{ij}$
(negative, from the Christoffel connection).}

\paragraph{(ii) Combining:} The $(0,0)$ and $(i,j)$ components of
$H_{\mu\nu}$ on FLRW are (using
$-2g_{00}\Box R = +2\Box R$,
$-\tfrac{1}{2}g_{00}R^2 = +\tfrac{1}{2}R^2$):
\begin{align}
  H_{00} &= 2\ddot{R} + 2\Box R + \tfrac{1}{2}R^2
  + 2R\cdot R_{00} \notag\\
  &= 2\ddot{R} + 2(-\ddot{R} - 3H\dot{R}) + \tfrac{1}{2}R^2
  - 6R(\dot{H} + H^2) \notag\\
  &= -6H\dot{R} + \tfrac{1}{2}R^2
  - 6R(\dot{H} + H^2),
  \label{eq:H00}\\[6pt]
  H_{ij} &= -2H\dot{R}\,a^2\delta_{ij}
  - 2\Box R\,a^2\delta_{ij}
  - \tfrac{1}{2}R^2\,a^2\delta_{ij}
  + 2R(\dot{H} + 3H^2)\,a^2\delta_{ij} \notag\\
  &= -2H\dot{R}\,a^2\delta_{ij}
  + 2(\ddot{R} + 3H\dot{R})\,a^2\delta_{ij}
  - \tfrac{1}{2}R^2\,a^2\delta_{ij}
  + 2R(\dot{H} + 3H^2)\,a^2\delta_{ij} \notag\\
  &= \bigl[2\ddot{R} + 4H\dot{R} + 2R(\dot{H} + 3H^2)
  - \tfrac{1}{2}R^2\bigr]\,a^2\delta_{ij}.
  \label{eq:Hij}
\end{align}

\signcorrection{LR audit correction: The original computation had
sign errors from $\nabla_0\nabla_0 R = -\ddot{R}$ and missing
$g_{00} = -1$ sign flips. With $\nabla_0\nabla_0 R = +\ddot{R}$
and correct metric sign factors, $H_{00}$ has $-6H\dot{R}$ and
$+\tfrac{1}{2}R^2$ (signs flipped from original), and $H_{ij}$ has
coefficient $4H\dot{R}$ (not $8H\dot{R}$). Verified by the trace
cross-check: $g^{\mu\nu}H_{\mu\nu} = -H_{00} + 3\cdot H_{ij}/(a^2)
= -6\Box R$, which holds with these corrected signs but fails with
the original ones.}

\begin{proposition}[De Sitter check]\label{prop:H-dS}
On de~Sitter ($\dot{H} = 0$, $R = 12H^2 = \mathrm{const}$):
$\dot{R} = \ddot{R} = 0$, so
\begin{align}
  H_{00}\big|_{\mathrm{dS}} &= +\tfrac{1}{2}(12H^2)^2
  - 6\cdot 12H^2\cdot H^2 = 72H^4 - 72H^4
  = 0, \notag\\
  H_{ij}\big|_{\mathrm{dS}} &= \bigl[2\cdot 12H^2 \cdot 3H^2
  - \tfrac{1}{2}(12H^2)^2\bigr]\,a^2\delta_{ij}
  = (72 - 72)\,H^4\,a^2\delta_{ij} = 0. \notag
\end{align}
\end{proposition}

\verified{De Sitter check: $H_{\mu\nu}|_{\mathrm{dS}} = 0$ directly
from the corrected formulas, consistent with NT-4b Check~5.
The $+\tfrac{1}{2}R^2$ and $2R\cdot R_{00} = -6R H^2$ terms cancel exactly.}

\subsection{The $\Theta^{(R)}_{\mu\nu}$ Terms on FLRW}

On FLRW, $R = R(t)$ implies $\nabla_\mu R = +\dot{R}\,\delta^0_\mu$,
and more generally $\nabla_\mu(\Box^k R) = +\frac{\dd}{\dd t}(\Box^k R)
\,\delta^0_\mu$. Therefore the alpha-insertion~\eqref{eq:Theta-R}
gives (with
$\nabla_\rho(\Box^k R)\,\nabla^\rho(\Box^{n-1-k} R)
= \nabla_0\cdot g^{00}\nabla_0 = -\dot{R}_k\,\dot{R}_{n-1-k}$):

\begin{align}
  \Theta^{(R)}_{00} &= -\sum_{n=1}^{\infty} c_n^{(2)}
  \sum_{k=0}^{n-1}
  \Bigl[\dot{R}_k\,\dot{R}_{n-1-k}
  - \tfrac{1}{2}\dot{R}_k\,\dot{R}_{n-1-k}\Bigr] \notag\\
  &= -\frac{1}{2}\sum_{n=1}^{\infty} c_n^{(2)}
  \sum_{k=0}^{n-1} \dot{R}_k\,\dot{R}_{n-1-k},
  \label{eq:Theta00}\\[6pt]
  \Theta^{(R)}_{ij} &= -\frac{1}{2}\,a^2\delta_{ij}
  \sum_{n=1}^{\infty} c_n^{(2)}
  \sum_{k=0}^{n-1} \dot{R}_k\,\dot{R}_{n-1-k},
  \label{eq:Thetaij}
\end{align}
where we abbreviated $R_k \equiv \Box^k R/\Lambda^{2k}$ and
$\dot{R}_k \equiv \frac{\dd}{\dd t}(\Box^k R/\Lambda^{2k})$.

\signcorrection{LR audit correction: (1) In $\Theta^{(R)}_{00}$,
the intermediate step uses ``$-$'' (not ``$+$'') for the second term
in the bracket, because $-(1/2)\,g_{00}\,\nabla_\rho\nabla^\rho
= -(1/2)(-1)(-\dot{R}_k\dot{R}_{n-1-k}) = -(1/2)\dot{R}_k\dot{R}_{n-1-k}$.
The final result $-(1/2)\mathcal{S}$ is unchanged.
(2) In $\Theta^{(R)}_{ij}$, the sign is $-1/2$ (not $+1/2$), because
$\nabla_i(\Box^k R) = 0$ and the only contribution is
$-(1/2)g_{ij}(-\dot{R}_k\dot{R}_{n-1-k}) = +(1/2)a^2\delta_{ij}
\dot{R}_k\dot{R}_{n-1-k}$ inside the bracket, which is then multiplied
by the overall $-\sum$ to give $-(1/2)a^2\delta_{ij}\mathcal{S}$.}

\begin{remark}[Structure of $\Theta^{(R)}$ on FLRW]
On FLRW, $\Theta^{(R)}_{\mu\nu}$ takes a ``perfect fluid'' form:
\begin{equation}\label{eq:Theta-pf}
  \Theta^{(R)}_{\mu\nu} = (\rho_\Theta + p_\Theta)\,u_\mu u_\nu
  + p_\Theta\,g_{\mu\nu},
\end{equation}
with effective energy density and pressure:
\begin{equation}\label{eq:rho-p-Theta}
  \rho_\Theta = -\frac{1}{2}\,\mathcal{S},
  \qquad
  p_\Theta = -\frac{1}{2}\,\mathcal{S},
  \qquad
  \mathcal{S} \equiv \sum_{n=1}^{\infty} c_n^{(2)}
  \sum_{k=0}^{n-1} \dot{R}_k\,\dot{R}_{n-1-k}.
\end{equation}
Note the stiff equation of state $p_\Theta = +\rho_\Theta$
($w = +1$), characteristic of kinetic-energy terms involving time
derivatives of $R$.
\end{remark}

\gap{} The exact relation between the infinite series $\mathcal{S}$
and the nonlocal operator $F_2(\Box/\Lambda^2)$ acting on $R$
needs to be made explicit. This connects to the integral
representation of the alpha-insertion.

%%======================================================================
\section{Modified Friedmann Equations}\label{sec:friedmann}
%%======================================================================

\subsection{Standard Friedmann Equations}

With a perfect fluid $T_{\mu\nu} = (\rho + p)\,u_\mu u_\nu
+ p\,g_{\mu\nu}$, the standard Friedmann equations are
\fromlit{Weinberg (2008), \S1.5}:
\begin{align}
  H^2 &= \frac{8\pi G}{3}\,\rho
  \quad \text{(first Friedmann)},
  \label{eq:fried1}\\
  \dot{H} + H^2 &= -\frac{4\pi G}{3}(\rho + 3p)
  \quad \text{(acceleration / Raychaudhuri)}.
  \label{eq:fried2}
\end{align}

\subsection{SCT-Modified Friedmann Equations}\label{sec:sct-friedmann}

On FLRW, the full field equations~\eqref{eq:full-eom} reduce to
(using $B_{\mu\nu} = 0$, $\Theta^{(C)} = 0$):
\begin{equation}\label{eq:eom-flrw}
  \frac{1}{\kappa^2}\,G_{\mu\nu}
  + \frac{1}{16\pi^2}\bigl[F_2(\Box/\Lambda^2,\xi)\,H_{\mu\nu}
  + \Theta^{(R)}_{\mu\nu}\bigr]
  = \frac{1}{2}\,T_{\mu\nu}.
\end{equation}

Taking the $(0,0)$ component and the trace, we obtain the modified
Friedmann equations.

\begin{theorem}[Modified first Friedmann equation]\label{thm:fried1-mod}
The $(0,0)$ component of~\eqref{eq:eom-flrw} gives:
\begin{equation}\label{eq:fried1-mod}
  3H^2 + \frac{\kappa^2}{16\pi^2}\Bigl[
  F_2(\Box/\Lambda^2,\xi)\,H_{00}
  + \Theta^{(R)}_{00}\Bigr]
  = \kappa^2\,\rho,
\end{equation}
where $H_{00}$ is given by~\eqref{eq:H00} and $\Theta^{(R)}_{00}$
by~\eqref{eq:Theta00}.
\end{theorem}

\begin{theorem}[Modified trace equation]\label{thm:trace-mod}
Taking the trace of~\eqref{eq:eom-flrw}:
\begin{equation}\label{eq:trace-mod}
  -\frac{R}{\kappa^2}
  + \frac{1}{16\pi^2}\bigl[
  -6\,F_2(\Box/\Lambda^2,\xi)\,\Box R
  + g^{\mu\nu}\Theta^{(R)}_{\mu\nu}\bigr]
  = \frac{1}{2}\,(-\rho + 3p),
\end{equation}
using $g^{\mu\nu}G_{\mu\nu} = -R$,
$g^{\mu\nu}H_{\mu\nu} = -6\Box R$ (NT-4b Prop.~3.1),
and $g^{\mu\nu}\Theta^{(R)}_{\mu\nu} = -\rho_\Theta + 3p_\Theta
= 2\rho_\Theta = -\mathcal{S}$ from~\eqref{eq:rho-p-Theta}
(using $p_\Theta = \rho_\Theta$, $w = +1$).
\end{theorem}

\subsection{Local Limit: $R + R^2$ Gravity (Starobinsky)}

In the local limit $F_2 \to F_2(0,\xi) = \alpha_R(\xi)/(16\pi^2)$,
$\Theta^{(R)} = 0$, the trace equation reduces to:
\begin{equation}\label{eq:trace-local}
  R + \frac{6\alpha_R(\xi)}{16\pi^2\kappa^{-2}}\,\Box R
  = -\kappa^2\,(\rho - 3p)/2.
\end{equation}

Comparing with Starobinsky's $R + R^2/(6M^2)$ model
\fromlit{Starobinsky (1980); Mukhanov \& Chibisov (1981)}: the
correspondence is $M^2 = 16\pi^2/(6\alpha_R(\xi)\kappa^2)
= 1/(6\alpha_R(\xi)\cdot 8\pi G)$.

\subsection{De Sitter Self-Consistency}

On de~Sitter ($H = \mathrm{const}$, $R = 12H^2$):
\begin{itemize}[nosep]
  \item $\dot{R} = 0$, $\Box R = 0$, $\Box^n R = 0$ for all
    $n \geq 1$.
  \item $H_{\mu\nu}|_{\mathrm{dS}} = 0$ (NT-4b, Prop.~4.2).
  \item $\Theta^{(R)}_{\mu\nu}|_{\mathrm{dS}} = 0$ (all terms
    involve $\nabla R$).
\end{itemize}
Therefore the modified Friedmann equation~\eqref{eq:fried1-mod}
reduces to $3H^2 = \kappa^2\rho$, i.e., standard Einstein gravity.
De~Sitter is an exact vacuum solution for $\rho = p = 0$ with
$H^2 = \Lambda_{\mathrm{eff}}/3$.

\verified{This is consistent with NT-4b Check~5 (de~Sitter exact
solution). The FLRW reduction inherits this check automatically.}

%%======================================================================
\section{Comparison with Nonlocal Cosmological Models}
\label{sec:comparison}
%%======================================================================

\subsection{Deser--Woodard Model}\label{sec:DW}

The DW model \cite{DW07,DW13,DW19} modifies gravity via:
\begin{equation}\label{eq:DW-action}
  S_{\mathrm{DW}} = \frac{1}{16\pi G}\int\dd^4 x\,\sqrt{-g}\,
  \Bigl[R + R\,f\Bigl(\frac{R}{\Box}\Bigr)\Bigr],
\end{equation}
where $f$ is a dimensionless function. The model introduces no
additional mass scale (both $R$ and $\Box$ have dimension
$[\mathrm{length}]^{-2}$).

\paragraph{Differences from SCT:}
\begin{enumerate}[nosep]
  \item DW uses $\Box^{-1}$ (retarded Green's function), which is
    infrared-motivated. SCT uses entire form factors
    $F_i(\Box/\Lambda^2)$, which are UV-motivated.
  \item DW modifies gravity in the IR (dark energy). SCT modifies
    it at the spectral cutoff scale $\Lambda$.
  \item DW has $c_T = c$ (GWs travel at the speed of light)
    \cite{NLA18}. SCT needs verification (see \S\ref{sec:gw-speed}).
\end{enumerate}

\paragraph{Key result \fromlit{Deser--Woodard (2013), \S2}:}
On FLRW, the DW model introduces an auxiliary field
$X \equiv \Box^{-1}R$ satisfying $\Box X = R$, which localizes the
dynamics. The modified Friedmann equations involve $X(t)$ and its
derivatives, producing an effective dark energy component.

\subsection{Maggiore--Mancarella RR Model}\label{sec:MM}

The RR model \cite{BDFM18,DFKKM14} adds $m^2\,R\,\Box^{-2}\,R$ to
the Einstein--Hilbert action:
\begin{equation}\label{eq:RR-action}
  S_{\mathrm{RR}} = \frac{1}{16\pi G}\int\dd^4 x\,\sqrt{-g}\,
  \Bigl[R - \frac{m^2}{6}\,R\,\frac{1}{\Box^2}\,R\Bigr].
\end{equation}

\paragraph{Key results \fromlit{Belgacem et al.\ (2018), \S3}:}
\begin{itemize}[nosep]
  \item The model introduces two auxiliary fields $U = -\Box^{-1}R$
    and $S = -\Box^{-1}U$ to localize the dynamics.
  \item On FLRW, the modified Friedmann equation becomes:
    \begin{equation}\label{eq:RR-friedmann}
      H^2 = \frac{8\pi G}{3}\rho
      + \frac{m^2}{6}\bigl[\text{terms involving }U, S, \dot{U},
      \dot{S}\bigr].
    \end{equation}
  \item The dark energy equation of state is phantom:
    $w_{\mathrm{DE}}(0) \approx -1.14$.
  \item GWs propagate at the speed of light: $c_T = c$.
  \item The model reduces $H_0$ tension to $2.0\sigma$.
\end{itemize}

\paragraph{Differences from SCT:}
The RR model uses $\Box^{-2}$ (IR nonlocality), while SCT uses
entire functions (UV nonlocality). The RR model has one free parameter
$m$ replacing $\Omega_\Lambda$; SCT has $\Lambda$ (spectral cutoff)
and $\xi$ (Higgs non-minimal coupling).

\subsection{Koshelev--Kumar--Starobinsky Nonlocal Inflation}
\label{sec:KKS}

The KKS approach \cite{KKS23} considers the action:
\begin{equation}\label{eq:KKS-action}
  S_{\mathrm{KKS}} = \frac{1}{2}\int\dd^4 x\,\sqrt{-g}\,
  \Bigl[M_p^2\,R + R\,\mathcal{F}_R(\Box_s)\,R
  + W_{\mu\nu\rho\sigma}\,\mathcal{F}_C(\Box_s)\,
  W^{\mu\nu\rho\sigma}\Bigr],
\end{equation}
where $\Box_s = \Box/\mathcal{M}_s^2$ and $\mathcal{F}_R$,
$\mathcal{F}_C$ are analytic infinite-derivative (AID) operators.

\paragraph{Key results \fromlit{KKS (2023), \S2--3}:}
\begin{itemize}[nosep]
  \item On FLRW, the Weyl sector drops out (same as SCT).
  \item The trace equation reduces to $\Box R = M^2 R$, the
    Starobinsky scalaron equation, under the condition
    $\mathcal{F}_R(M^2/\mathcal{M}_s^2) = M_p^2/(6M^2)$.
  \item Ghost-free form factors require
    $\mathcal{O}_0 = (1 - \Box/M^2)\,e^{H_0(\Box_s)}$ and
    $\mathcal{O}_2 = e^{H_2(\Box_s)}$ where $H_0$, $H_2$ are
    entire functions.
  \item The form factors are \emph{background-dependent} on curved
    spacetimes, necessitating cubic and higher curvature terms
    for a background-independent formulation.
  \item On FLRW, the scalaron equation of motion provides
    $R^2$-like inflation.
\end{itemize}

\paragraph{Relevance to SCT:} The SCT form factors $F_1$, $F_2$
are derived from the spectral action, not postulated. The KKS
ghost-freedom conditions correspond to the SCT requirement that
$F_1$, $F_2$ are entire (NT-2). The scalaron mass in KKS corresponds
to the SCT effective mass $m_0(\xi) = \Lambda/\sqrt{6(\xi - 1/6)^2}$.

\subsection{Biswas--Koshelev--Mazumdar Bouncing Cosmology}
\label{sec:BKM}

The BKM approach \cite{BMS06,BKM10} obtains bouncing cosmological
solutions from ghost-free nonlocal gravity:
\begin{equation}\label{eq:BKM-action}
  S_{\mathrm{BKM}} = \frac{M_p^2}{2}\int\dd^4 x\,\sqrt{-g}\,
  \bigl[R + R\,\mathcal{F}(\Box)\,R\bigr].
\end{equation}

\paragraph{Key results \fromlit{BKM (2010), \S3}:}
\begin{itemize}[nosep]
  \item Stable inflationary attractor solutions exist with positive
    cosmological constant.
  \item Cyclic universe solutions emerge with negative cosmological
    constant.
  \item Super-Hubble perturbations are tracked through the bounce.
  \item Bouncing solutions are free from perturbative instability.
\end{itemize}

\paragraph{Relevance to SCT:} The BKM action is structurally identical
to the $R^2$ sector of the SCT spectral action, with
$\mathcal{F}(\Box) \leftrightarrow F_2(\Box/\Lambda^2,\xi)/(16\pi^2)$.
The SCT form factor $F_2$ is derived from first principles (spectral
action + heat kernel), while BKM postulates $\mathcal{F}$.

%%======================================================================
\section{Gravitational Wave Speed}\label{sec:gw-speed}
%%======================================================================

\subsection{The GW170817 Constraint}

The joint detection of GW170817 and GRB~170817A established
\fromlit{LIGO--Virgo Collaboration (2017)}:
\begin{equation}\label{eq:gw170817}
  \bigl|c_T/c - 1\bigr| \leq 1 \times 10^{-15},
\end{equation}
where $c_T$ is the speed of gravitational waves. This constraint
applies at $z \approx 0$ (the source was at 40~Mpc).

\subsection{Tensor Perturbations on FLRW}\label{sec:tensor-pert}

Consider tensor perturbations $h_{ij}^{\mathrm{TT}}$ on the FLRW
background \fromlit{Mukhanov (2005), \S7; Maggiore (2008), \S19}:
\begin{equation}\label{eq:metric-tensor-pert}
  \dd s^2 = -\dd t^2 + a^2(t)\bigl(\delta_{ij}
  + h_{ij}^{\mathrm{TT}}\bigr)\,\dd x^i\,\dd x^j,
\end{equation}
where $h_{ij}^{\mathrm{TT}}$ is transverse and traceless:
$\partial_i h^{ij}_{\mathrm{TT}} = 0$, $h^i{}_i = 0$.

In GR, the tensor perturbation equation is:
\begin{equation}\label{eq:gw-GR}
  \ddot{h}_{ij} + 3H\dot{h}_{ij} - \frac{\nabla^2}{a^2}\,h_{ij} = 0.
\end{equation}

\subsection{GW Propagation in SCT on FLRW}

\begin{tcolorbox}[colback=yellow!5!white, colframe=orange!60!black,
  title={\textbf{Critical: The Weyl Sector Contributes to
  Perturbations}}]
Although $C_{\mu\nu\rho\sigma} = 0$ on the FLRW \emph{background},
the perturbation $\delta C_{\mu\nu\rho\sigma} \neq 0$ for tensor
modes. Therefore the $F_1$ (Weyl) sector \emph{does} contribute to
the tensor perturbation equation, even though it does not appear in
the background Friedmann equations.

This is a standard feature of higher-derivative gravity theories:
the linearized equation of motion for perturbations around a curved
background can involve terms from sectors that vanish on the
background itself.
\end{tcolorbox}

From the NT-4b field equations linearized on FLRW, the tensor
perturbation equation takes the schematic form:
\begin{equation}\label{eq:gw-sct}
  \bigl[1 + \mathcal{G}_T(\Box/\Lambda^2)\bigr]\,
  \bigl(\ddot{h}_{ij} + 3H\dot{h}_{ij}
  - \frac{\nabla^2}{a^2}\,h_{ij}\bigr) = 0,
\end{equation}
where $\mathcal{G}_T$ encodes the nonlocal modification from
the Weyl sector.

\gap{} The exact form of $\mathcal{G}_T$ on a general FLRW
background needs to be derived. On flat background, it reduces to
the NT-4a result $\Pi_{\mathrm{TT}}(z) = 1 + (13/60)\,z\,\hat{F}_1(z)$.

\paragraph{The GW speed:}
For a mode $h_{ij} \propto e^{i(\omega t - \mathbf{k}\cdot\mathbf{x})}$,
the dispersion relation from~\eqref{eq:gw-sct} gives:
\begin{equation}\label{eq:gw-dispersion}
  \omega^2 = \frac{k^2}{a^2}\cdot c_T^2(k,t),
\end{equation}
where $c_T$ depends on $k$ in general (nonlocal dispersion).

\paragraph{Low-$k$ limit:}
In the limit $k \to 0$ (long wavelengths, relevant for GW170817):
\begin{equation}\label{eq:cT-lowk}
  c_T^2(k \to 0) = 1,
\end{equation}
because the form factors reduce to their $z = 0$ values and the
effective action becomes the local Einstein--Hilbert action
(plus local curvature-squared terms that do not modify the GW speed
in $d = 4$).

\begin{proposition}[GW speed in SCT at low momenta]
\label{prop:cT}
At momenta $k \ll \Lambda$ (equivalently $z = k^2/\Lambda^2 \ll 1$),
the SCT prediction for the GW speed on FLRW is
$c_T = c$, consistent with the GW170817 bound~\eqref{eq:gw170817}.
\end{proposition}

This follows from the fact that:
\begin{enumerate}[nosep]
  \item In the local limit ($z \to 0$), the Weyl sector gives
    $2F_1(0)\,B_{\mu\nu} = (13/60)\cdot B_{\mu\nu}/(16\pi^2)$.
  \item The local $C^2$ term does not modify the GW speed in 4D
    \fromlit{Stelle (1977); Nersisyan et al.\ (2018), \S3}.
  \item The $R^2$ term does not affect tensor perturbations at
    leading order (it couples to the scalar sector only).
\end{enumerate}

\paragraph{Comparison:} Both the DW model \cite{NLA18} and the
Maggiore RR model \cite{BDFM18} have $c_T = c$ exactly. The
Horndeski models with $G_4(X) \neq 0$ or $G_5 \neq 0$ are
excluded \cite{NLA18}. The SCT prediction $c_T(k \to 0) = c$
places it in the same class as DW and RR.

\gap{} The high-$k$ behavior of $c_T(k)$ (UV dispersion) is
determined by $F_1(z)$ and may show deviations from $c$ at
$k \sim \Lambda$. This is expected from the MR-3 causality
analysis (micro-causality violated at $\ell \sim 1/\Lambda$).

%%======================================================================
\section{Cosmological Perturbation Theory on FLRW}
\label{sec:perturbations}
%%======================================================================

\subsection{SVT Decomposition}

The perturbed FLRW metric in the Newtonian (longitudinal) gauge is
\fromlit{Mukhanov (2005), \S7; Bardeen (1980)}:
\begin{equation}\label{eq:perturbed-metric}
  \dd s^2 = -(1 + 2\Phi)\,\dd t^2
  + a^2(t)\bigl[(1 - 2\Psi)\,\delta_{ij}
  + h_{ij}^{\mathrm{TT}}\bigr]\,\dd x^i\,\dd x^j,
\end{equation}
where $\Phi$ and $\Psi$ are the Bardeen scalar potentials and
$h_{ij}^{\mathrm{TT}}$ is the transverse-traceless tensor perturbation.
Vector perturbations decay in an expanding universe and are omitted.

\subsection{Scalar Perturbations and Gravitational Slip}

In GR, $\Phi = \Psi$ in the absence of anisotropic stress. In
modified gravity, the gravitational slip parameter
\fromlit{Amendola et al.\ (2008)}:
\begin{equation}\label{eq:grav-slip}
  \eta \equiv -\frac{\Phi}{\Psi}
\end{equation}
can differ from unity. For the DW model,
$\eta \approx -0.05 + 1.79\,a^2 - 0.54\,a^3$ \cite{NLA18}.

In SCT on FLRW, the $R^2$ sector generates anisotropic stress via
the $H_{\mu\nu}$ tensor, producing $\Phi \neq \Psi$. The Weyl
sector $F_1$ also contributes to scalar perturbations through
$\delta C_{\mu\nu\rho\sigma}$.

\gap{} The explicit expressions for $\eta(k, z)$ in SCT need
to be derived in the NT-4c derivation phase. This is important
for INF-1 ($n_s$, $r$) and MT-2 ($H_0$, $S_8$ tensions).

\subsection{Inflationary Observables (Preview for INF-1)}

For $R^2$-like inflation driven by the $F_2$ sector, the key
observables are \fromlit{KKS (2023), \S4}:
\begin{align}
  n_s &= 1 - \frac{2}{N}
  \quad \text{(spectral index, at leading order)},
  \label{eq:ns}\\
  r &= \frac{12}{N^2}
  \quad \text{(tensor-to-scalar ratio, at leading order)},
  \label{eq:r}
\end{align}
where $N$ is the number of e-folds. For $N = 55$:
$n_s \approx 0.964$, $r \approx 0.004$, consistent with Planck 2018
constraints $n_s = 0.9649 \pm 0.0042$ and $r < 0.064$ (95\% CL).

In the SCT framework, nonlocal corrections from $F_2(\Box/\Lambda^2)$
modify both $n_s$ and $r$ at sub-leading order. These corrections
depend on the ratio $M^2/\mathcal{M}_s^2$ (scalaron mass vs.\
nonlocality scale).

\gap{} The computation of $n_s$ and $r$ in SCT, including
nonlocal corrections, is deferred to INF-1.

%%======================================================================
\section{Notation Table}\label{sec:notation}
%%======================================================================

\begin{center}
\begin{longtable}{ll}
\toprule
\textbf{Symbol} & \textbf{Definition / Convention} \\
\midrule
\endfirsthead
\toprule
\textbf{Symbol} & \textbf{Definition / Convention} \\
\midrule
\endhead
\bottomrule
\endfoot
$g_{\mu\nu}$ & Lorentzian metric, signature $(-,+,+,+)$ \\
$a(t)$ & Scale factor \\
$H = \dot{a}/a$ & Hubble parameter \\
$\mathcal{H} = a'/a$ & Conformal Hubble parameter \\
$\tau$ & Conformal time: $\dd\tau = \dd t/a$ \\
$R$ & Ricci scalar: $6(\dot{H} + 2H^2)$ on FLRW \\
$\Box$ & d'Alembertian: $-\ddot{\phantom{x}} - 3H\dot{\phantom{x}}$
  for scalars on FLRW \\
$\Lambda$ & Spectral cutoff scale \\
$\xi$ & Higgs non-minimal coupling \\
$\kappa^2$ & $8\pi G$ \\
$\alpha_C$ & $13/120$ (Weyl coefficient) \\
$\alpha_R(\xi)$ & $2(\xi - 1/6)^2$ (Ricci-scalar coefficient) \\
$F_1(z)$ & Weyl form factor, $z = \Box/\Lambda^2$ \\
$F_2(z,\xi)$ & $R^2$ form factor \\
$B_{\mu\nu}$ & Bach tensor (from $\delta C^2/\delta g^{\mu\nu}$) \\
$H_{\mu\nu}$ & $R^2$-sector tensor \\
$\Theta^{(R)}_{\mu\nu}$ & BV alpha-insertion, $R^2$ sector \\
$\Theta^{(C)}_{\mu\nu}$ & BV alpha-insertion, Weyl sector \\
$\Phi, \Psi$ & Bardeen scalar potentials \\
$h_{ij}^{\mathrm{TT}}$ & Tensor perturbation \\
$c_T$ & Gravitational wave speed \\
$n_s$ & Scalar spectral index \\
$r$ & Tensor-to-scalar ratio \\
$N$ & Number of e-folds \\
\bottomrule
\end{longtable}
\end{center}

%%======================================================================
\section{Summary of Gaps for the Derivation Phase}\label{sec:gaps}
%%======================================================================

\begin{center}
\begin{tabular}{clp{8cm}l}
\toprule
\textbf{\#} & \textbf{Gap} & \textbf{Description} & \textbf{Priority} \\
\midrule
G1 & $H_{\mu\nu}$ on FLRW &
  Verify the explicit $(0,0)$ and $(i,j)$ components of $H_{\mu\nu}$
  on FLRW, including the action of $F_2(\Box/\Lambda^2)$.
  & HIGH \\
G2 & $\Theta^{(R)}$ convergence &
  Prove convergence of the alpha-insertion series on FLRW.
  Establish connection between the series $\mathcal{S}$ and the
  integral representation.
  & HIGH \\
G3 & Tensor perturbation equation &
  Derive the exact form of $\mathcal{G}_T(\Box/\Lambda^2)$ on FLRW
  from the linearization of the Weyl sector around the FLRW
  background. Extract $c_T(k)$.
  & HIGH \\
G4 & Scalar perturbation equations &
  Derive the coupled equations for $\Phi$, $\Psi$ from the
  linearization of both the $R^2$ and Weyl sectors on FLRW.
  Extract $\eta(k, z)$.
  & MEDIUM \\
G5 & Eigenvalue ansatz &
  Test the KKS ansatz $\Box R = M^2 R$ in the SCT context.
  Determine whether the SCT form factor $F_2$ admits this
  simplification on the FLRW background.
  & MEDIUM \\
G6 & Conformal coupling limit &
  At $\xi = 1/6$, $\alpha_R = 0$ and the $R^2$ sector vanishes.
  Verify that the tensor perturbation equation still receives
  contributions from the Weyl sector $F_1$ (which has $\alpha_C =
  13/120 \neq 0$).
  & HIGH \\
G7 & Energy conditions &
  Check whether the effective energy density and pressure from
  the nonlocal terms satisfy standard energy conditions (null,
  weak, strong).
  & LOW \\
\bottomrule
\end{tabular}
\end{center}

%%======================================================================
\section{Benchmarks for the Derivation Phase}\label{sec:benchmarks}
%%======================================================================

The following benchmarks must be satisfied by any correct derivation:

\begin{enumerate}[label=\textbf{B\arabic*}]
  \item \textbf{De Sitter:} The modified Friedmann equations must
    reduce to $3H^2 = \kappa^2\rho$ on de~Sitter, consistent with
    NT-4b Check~5.
  \item \textbf{Local limit:} Setting $F_2 \to F_2(0,\xi)$,
    $\Theta^{(R)} \to 0$, the equations must reproduce the local
    $R + \beta R^2$ Friedmann equations with
    $\beta = \alpha_R(\xi)/(16\pi^2)$.
  \item \textbf{Conformal coupling:} At $\xi = 1/6$,
    $\alpha_R = 0$, and the \emph{background} Friedmann equations
    must reduce to standard GR.
  \item \textbf{GW speed:} $c_T(k \to 0) = c$ at low momenta,
    consistent with GW170817.
  \item \textbf{Linearized recovery:} On flat background ($a = 1$,
    $H = 0$), the perturbation equations must reproduce NT-4a
    $\Pi_{\mathrm{TT}}$ and $\Pi_{\mathrm{s}}$.
  \item \textbf{Energy conservation:} The modified equations must be
    consistent with $\nabla^\mu T_{\mu\nu} = 0$ (Bianchi identity).
  \item \textbf{Trace consistency:} The trace of the modified
    Friedmann equations must give the Raychaudhuri equation with the
    same nonlocal corrections.
  \item \textbf{KKS scalaron equation:} Under the ansatz
    $\Box R = M^2 R$, the SCT trace equation must reduce to the
    Starobinsky scalaron equation.
\end{enumerate}

%%======================================================================
\begin{thebibliography}{99}
%%======================================================================

\bibitem{DW07}
S.~Deser and R.~P.~Woodard,
``Nonlocal cosmology,''
\emph{Phys.\ Rev.\ Lett.} \textbf{99} (2007) 111301,
\href{https://arxiv.org/abs/0706.2151}{arXiv:0706.2151}.

\bibitem{DW13}
S.~Deser and R.~P.~Woodard,
``Observational viability and stability of nonlocal cosmology,''
\emph{JCAP} \textbf{1311} (2013) 036,
\href{https://arxiv.org/abs/1307.6639}{arXiv:1307.6639}.

\bibitem{DW19}
S.~Deser and R.~P.~Woodard,
``Nonlocal cosmology~II --- Cosmic acceleration without fine tuning
or dark energy,''
\emph{JCAP} \textbf{1906} (2019) 034,
\href{https://arxiv.org/abs/1902.08075}{arXiv:1902.08075}.

\bibitem{BDFM18}
E.~Belgacem, Y.~Dirian, S.~Foffa, and M.~Maggiore,
``Nonlocal gravity. Conceptual aspects and cosmological predictions,''
\emph{JCAP} \textbf{1803} (2018) 002,
\href{https://arxiv.org/abs/1712.07066}{arXiv:1712.07066}.

\bibitem{DFKKM14}
Y.~Dirian, S.~Foffa, N.~Khosravi, M.~Kunz, and M.~Maggiore,
``Cosmological perturbations and structure formation in nonlocal
infrared modifications of general relativity,''
\emph{JCAP} \textbf{1406} (2014) 033,
\href{https://arxiv.org/abs/1403.6068}{arXiv:1403.6068}.

\bibitem{BDFM18gw}
E.~Belgacem, Y.~Dirian, S.~Foffa, and M.~Maggiore,
``Modified gravitational-wave propagation and standard sirens,''
\emph{Phys.\ Rev.\ D} \textbf{98} (2018) 023510,
\href{https://arxiv.org/abs/1805.08731}{arXiv:1805.08731}.

\bibitem{KKS23}
A.~S.~Koshelev, K.~Sravan~Kumar, and A.~A.~Starobinsky,
``Cosmology in nonlocal gravity,''
in \emph{Handbook of Quantum Gravity}, Springer, 2023,
\href{https://arxiv.org/abs/2305.18716}{arXiv:2305.18716}.

\bibitem{BMS06}
T.~Biswas, A.~Mazumdar, and W.~Siegel,
``Bouncing universes in string-inspired gravity,''
\emph{JCAP} \textbf{0603} (2006) 009,
\href{https://arxiv.org/abs/hep-th/0508194}{hep-th/0508194}.

\bibitem{BKM10}
T.~Biswas, T.~Koivisto, and A.~Mazumdar,
``Towards a resolution of the cosmological singularity in non-local
higher derivative theories of gravity,''
\emph{JCAP} \textbf{1011} (2010) 008,
\href{https://arxiv.org/abs/1005.0590}{arXiv:1005.0590}.

\bibitem{BCM19}
F.~Briscese, G.~Calcagni, and L.~Modesto,
``Nonlinear stability in nonlocal gravity,''
\emph{Phys.\ Rev.\ D} \textbf{99} (2019) 084041,
\href{https://arxiv.org/abs/1901.03267}{arXiv:1901.03267}.

\bibitem{CM22}
A.~Bas~i~Beneito, G.~Calcagni, and L.~Rachwa\l{},
``Classical and quantum nonlocal gravity,''
in \emph{Handbook of Quantum Gravity}, Springer, 2022,
\href{https://arxiv.org/abs/2211.05606}{arXiv:2211.05606}.

\bibitem{NLA18}
H.~Nersisyan, N.~A.~Lima, and L.~Amendola,
``Gravitational wave speed: Implications for models without a
mass scale,''
\emph{Phys.\ Lett.\ B} \textbf{783} (2018) 188--193,
\href{https://arxiv.org/abs/1801.06683}{arXiv:1801.06683}.

\bibitem{NO08}
S.~Nojiri and S.~D.~Odintsov,
``Modified gravity as realistic candidate for dark energy, inflation
and dark matter,''
\href{https://arxiv.org/abs/0810.1557}{arXiv:0810.1557}.

\bibitem{Ste77}
K.~S.~Stelle,
``Renormalization of higher-derivative quantum gravity,''
\emph{Phys.\ Rev.\ D} \textbf{16} (1977) 953--969.

\bibitem{Ste78}
K.~S.~Stelle,
``Classical gravity with higher derivatives,''
\emph{Gen.\ Rel.\ Grav.} \textbf{9} (1978) 353--371.

\bibitem{BV87}
A.~O.~Barvinsky and G.~A.~Vilkovisky,
``The generalized Schwinger-DeWitt technique in gauge theories and
quantum gravity,''
\emph{Phys.\ Rep.} \textbf{119} (1985) 1--74;
\emph{Nucl.\ Phys.\ B} \textbf{282} (1987) 163--188.

\bibitem{BV90}
A.~O.~Barvinsky and G.~A.~Vilkovisky,
``Covariant perturbation theory (II): second order in the curvature,''
\emph{Nucl.\ Phys.\ B} \textbf{333} (1990) 471--511.

\bibitem{Star80}
A.~A.~Starobinsky,
``A new type of isotropic cosmological models without singularity,''
\emph{Phys.\ Lett.\ B} \textbf{91} (1980) 99--102.

\bibitem{CPR09}
A.~Codello, R.~Percacci, and C.~Rahmede,
``Investigating the ultraviolet properties of gravity with a
Wilsonian renormalization group equation,''
\emph{Ann.\ Phys.} \textbf{324} (2009) 414--469,
\href{https://arxiv.org/abs/0805.2909}{arXiv:0805.2909}.

\bibitem{CZ12}
A.~Codello and O.~Zanusso,
``On the non-local heat kernel expansion,''
\emph{J.\ Math.\ Phys.} \textbf{54} (2013) 013513,
\href{https://arxiv.org/abs/1203.2034}{arXiv:1203.2034}.

\bibitem{LIGO17}
B.~P.~Abbott \emph{et al.}\ (LIGO Scientific Collaboration and
Virgo Collaboration),
``Gravitational waves and gamma-rays from a binary neutron star
merger: GW170817 and GRB~170817A,''
\emph{Astrophys.\ J.\ Lett.} \textbf{848} (2017) L13.

\bibitem{Muk05}
V.~F.~Mukhanov,
\emph{Physical Foundations of Cosmology},
Cambridge University Press, 2005.

\bibitem{Wei08}
S.~Weinberg,
\emph{Cosmology},
Oxford University Press, 2008.

\bibitem{Car04}
S.~M.~Carroll,
\emph{Spacetime and Geometry: An Introduction to General Relativity},
Addison-Wesley, 2004.

\bibitem{Bar80}
J.~M.~Bardeen,
``Gauge-invariant cosmological perturbations,''
\emph{Phys.\ Rev.\ D} \textbf{22} (1980) 1882--1905.

\end{thebibliography}

\end{document}
